
% Authoritative LaTeX manuscript; Markdown is a generated compatibility projection.

An engineering question may span several models without requiring one larger stored model. Keep question-specific representations separate and join them through stable identities when a cross-model property must be evaluated.

\textbf{Engineering question.} Does an important user journey exercise its declared endpoints, receive the required test coverage and placement, and execute under the correct host policy?

\autoref{fig-g7-joined-scenario} traces the join.

\begin{figure}[htbp]
\centering
\includegraphics[width=\linewidth]{book/latex/assets/appendix-g-7-composition-by-reference-44bd4fbcec.pdf}
\caption{\emph{Composition by reference.} Stable identities join the user-journey, service, coverage, test-placement, and execution-policy views. The resulting joined view supports cross-model queries without duplicating those facts into a single stored model.}
\label{fig-g7-joined-scenario}
\end{figure}

\textbf{Property.} The joined view states cross-model properties that no single view can:

\begin{itemize}
\tightlist
\item
  \textbf{Every declared journey step has the required test coverage.}
\item
  \textbf{Declared journey dependencies agree with the actual call sites.}
\item
  \textbf{Declared endpoints are exercised.}
\item
  \textbf{Test placement reflects journey criticality.}
\item
  \textbf{Execution respects the selected host policy.}
\end{itemize}

\textbf{The reusable pattern.} State the engineering question, select the relevant views, join them on stable identities, and evaluate the cross-view property. The scenario is a traversal across representations, not an additional stored model.
